%
% konsequenzen.tex -- Konsequenzen
%
% (c) 2021 Prof Dr Andreas Müller, OST Ostschweizer Fachhochschule
%
\bgroup
\begin{frame}[t]
\setlength{\abovedisplayskip}{5pt}
\setlength{\belowdisplayskip}{5pt}
\frametitle{Konsequenzen der Holomorphie}
\begin{definition}[holomorph]
Eine Funktion $f\colon U\to\mathbb{C}$ heisst \emph{holomorph}, wenn sie
komplex differenzierbar ist:
\[
f'(z)
=
\frac{\partial f}{\partial z}
=
\frac12\biggl(
\frac{\partial f}{\partial x}
-
i
\frac{\partial f}{\partial y}
\biggr),
\qquad
\frac{\partial f}{\partial \bar{z}}
=
\frac{1}{2}\biggl(
\frac{\partial f}{\partial x}
+i
\frac{\partial f}{\partial y}
\biggr)
=
0.
\]
\end{definition}
%
\only<2-3>{
\begin{satz}[Cauchy-Riemann-Differentialgleichungen]
Realteil $u$ und Imaginärteil $v$ einer holomorphen Funktion
erfüllen die \emph{Cauchy-Riemann-Differentialgleichungen}
\[
\frac{\partial u}{\partial x} = \frac{\partial v}{\partial y}
\qquad\text{und}\qquad
\frac{\partial u}{\partial y} = -\frac{\partial v}{\partial x}
\]
\end{satz}
}
%
\only<3-5>{
\begin{satz}[harmonisch]
Realteil $u$ und Imaginärteil $v$ einer holomorphen Funktion sind
harmonisch $\Delta u=0$ und $\Delta v=0$.
\end{satz}
}
%
\only<4-5>{
\begin{satz}[Maximum-Minimum-Prinzip]
$u$ und $v$ erfüllen das Maximum-/Minimum-Prinzip: Maximum und
Minimum sind nicht im Inneren des Definitionsgebietes.
\end{satz}
}
%
\only<5>{
\begin{satz}[Mittelwerteigenschaft]
Ist $f$ holomorph, dann erfüllen $u$ und $v$ die Mittelwerteigenschaft.
\end{satz}
}
%
\only<6-7>{
\begin{satz}
Ist $\gamma$ ein geschlossener, zusammenziehbarer Weg im Definitionsgebiet
einer holomorphen Funktion $f$, dann ist
\[
\oint_\gamma f(z)\,dz = 0.
\]
\end{satz}
}
%
\only<7>{\vspace*{-0.7cm}}
\only<8->{\vspace*{-0.3cm}}
\only<7-9>{
\begin{satz} Ist $f$ eine holomorphe Funktion und $\gamma$ ein geschlossener
Weg, der den Punkt $a$ einmal im Gegenuhrzeigersinn umläuft, dann gilt
\[
f^{(k)}(a)
=
\frac{k!}{2\pi i}
\oint_\gamma \frac{f(z)}{(z-1)^{k+1}}\,dz.
\]
\end{satz}
}
%
\only<8-9>{
\vspace*{-0.8cm}
\begin{satz}
Ist $f$ holomorph, dann ist $f$ auch analytisch, d.~h.~hat eine konvergente
Potenzreihenentwicklung.
\end{satz}
}
%
\only<9>{
\vspace*{-0.1cm}
\begin{satz}
Holomorphe Funktionen sind beliebig oft stetig differenzierbar, d.~h.~glatt.
\end{satz}
}

\end{frame}
\egroup
